\begin{textsample}{POS Dim 1 – pi – Score 36.00 – pi/pi\_em\_43.txt}  \label{ex:f1_pos_090}
Competências e Habilidades com a Competência Geral 2 da BNCC, no \textbf{que} se refere ao exercício da curiosidade intelectual \textbf{que} utiliza o conhecimento para investigar, refletir e criar soluções em diferentes situações. ( CE02 ) Articular conhecimentos matemáticos ao propor e/ou participar de ações para investigar desafios do mundo \textbf{contemporâneo} e tomar decisões éticas e socialmente responsáveis, com base na análise de problemas de urgência social, como os voltados a situações de saúde, sustentabilidade, das implicações da tecnologia no mundo do trabalho, entre outros, recorrendo a conceitos, procedimentos e linguagens próprios da Matemática.
Esta competência está voltada para a investigação de ações possíveis de serem realizadas a partir de análises com base em elementos da matemática, tendo em vista aspectos éticos e socialmente responsáveis. \textbf{Note} \textbf{que} a ênfase dessa competência está na intervenção na realidade.
Essas ações de intervenção poderão fornecer \textbf{um} excelente contexto para aprender novos conceitos matemáticos ao mesmo tempo em \textbf{que} permitem aos alunos vivenciar a \textbf{aplicabilidade} da matemática.
A competência 2 coloca o estudante como personagem atuante em sua comunidade local e no mundo \textbf{globalizado}.
As ações de propor e participar fazem referência à capacidade de ser parte de algo, compartilhar saberes com o outro e colaborar conjuntamente para a produção de algo.
Destaca -se também o papel da investigação por parte do estudante, o \textbf{que} pressupõe a observação dos desafios presentes em sua comunidade local/global, a elaboração de hipóteses \textbf{que} as descrevam, o tratamento dos dados associados à situação envolvida, a análise dos resultados obtidos e, por fim, a tomada de decisão a partir das conclusões obtidas.
Ao desenvolver essa competência, pode -se afirmar \textbf{que} o estudante avança em relação ao entendimento de \textbf{que} os Projetos de Vida não são apenas no âmbito profissional, mas também nas dimensões pessoal e social/cidadã.
Está associada à Competência Geral 7 da BNCC \textbf{que} permite ao estudante \textbf{um} \textbf{posicionamento} ético em relação ao cuidado de si mesmo, dos outros e do planeta, bem como com a Competência Geral 10 \textbf{que} trata do agir pessoal e coletivo com tomada de decisões, a partir de princípios éticos, democráticos, inclusivos, sustentáveis e solidários. ( CE03 ) Utilizar estratégias, conceitos e procedimentos matemáticos, em seus campos – Aritmética, Álgebra, Grandezas e Medidas, Geometria, Probabilidade e Estatística – para interpretar, construir modelos e resolver problemas em diversos contextos, analisando a plausibilidade dos resultados e a adequação das soluções propostas, de modo a construir argumentação consistente.
A C3 prevê a utilização de estratégias, conceitos, definições e procedimentos matemáticos para construir modelos e resolver problemas. \textbf{Uma} vez construído o modelo e resolvido o problema, é \textbf{preciso} se debruçar sobre essa solução de modo a analisar a adequação e veracidade do elaborado para convencer a si, aos colegas e ao professor de \textbf{que} sua conclusão é correta. \textbf{Convém} também observar \textbf{que} no processo de resolução de problemas o estudante deve ser motivado a : questionar … formular … testar e validar hipóteses … buscar contraexemplos … modelar situações … verificar adequação da resposta.
E, como consequência, desenvolver linguagens e construir formas de pensar \textbf{que} o levem a refletir e agir de maneira crítica.
É importante \textbf{frisar} \textbf{que} a referida competência não se \textbf{restringe} apenas à resolução de problemas, mas também trata de sua elaboração.
Isso \textbf{revela} \textbf{uma} concepção da resolução de problemas além da \textbf{mera} aplicação de \textbf{um} conjunto de regras.
Outro grande destaque refere -se à modelagem matemática como a construção de modelos matemáticos \textbf{que} sirvam para generalizar ideias ou para descrever situações semelhantes.
Essa competência tem estreita relação com a Competência Geral 2 da BNCC, no sentido da capacidade de formular e resolver problemas, e com a Competência Geral 4, \textbf{que} reforça a importância de saber utilizar as diferentes linguagens para expressar ideias e informações para a comunicação mútua. ( CE04 ) Compreender e utilizar, com flexibilidade e fluidez, diferentes registros de representação matemáticos ( algébrico, geométrico, estatístico, computacional, entre outros ), na busca de solução e comunicação de resultados de problemas, de modo a favorecer a construção e o desenvolvimento do \textbf{raciocínio} matemático.
Nesta competência há a proposição do uso flexível dos diferentes registros de representação da linguagem matemática.
O domínio de tais registros amplia a capacidade de análise das situações representadas e \textbf{aprimora} tanto os aspectos de formulação como a comunicação de resultados.
As aprendizagens previstas nessa competência contribuem ainda para melhorar a capacidade de \textbf{argumentar} e identificar \textbf{raciocínios} falsos.
A competência 4 complementa as demais no sentido de \textbf{que} utilizar, interpretar e resolver situações-problema se faz pela comunicação das ideias dos estudantes por meio da linguagem matemática.
Transitar entre os diversos tipos de representações ( simbólica, algébrica, gráfica, textual etc. ) permite a compreensão mais profunda dos conceitos e ideias da matemática.
A representação de \textbf{uma} mesma situação de diferentes formas estabelece conexões \textbf{que} possibilitam resolver problemas matemáticos usando estratégias diversas.
Além disso, a capacidade de elaborar modelos matemáticos para expressar situações \textbf{implica} e \textbf{revela} a aprendizagem, além de potencializar o letramento matemático.
Essa competência está relacionada ao desenvolvimento das Competências Gerais 4 e 5 da BNCC, \textbf{uma} vez \textbf{que} a linguagem utilizada de modo flexível permite expressar ideias e informações \textbf{que} facilitam o entendimento e ampliar o repertório de formas de expressão, inclusive a digital com espaço para autoria pessoal e criatividade do estudante. ( CE05 ) Investigar e estabelecer conjecturas a respeito de diferentes conceitos e propriedades matemáticas, empregando recursos e estratégias como observação de padrões, experimentações e tecnologias digitais, identificando a necessidade, ou não, de \textbf{uma} demonstração cada vez mais formal na validação das referidas conjecturas.
A competência 5 prevê \textbf{que} os alunos potencializem seus processos de abstração e generalização na elaboração de conjecturas.
Busquem recursos \textbf{que} permitam \textbf{uma} validação mais formal da conjectura elaborada ou sua refutação.
Os alunos devem utilizar o \textbf{raciocínio} indutivo \textbf{baseado} na observação de padrões e em regularidades.
No entanto, eles deverão reconhecer \textbf{que} o \textbf{raciocínio} indutivo e as verificações meramente empíricas não são apropriadas para validar propriedades, podendo levar, algumas vezes, a conclusões equivocadas.
Desse modo, é necessário reconhecerem \textbf{que} as conjecturas formuladas devem ser justificadas por meio de resultados matemáticos já demonstrados.
Assim, a competência 5 tem como objetivo principal \textbf{que} os estudantes se apropriem da forma de pensar matemática, como ciência com \textbf{uma} forma específica de validar suas conclusões pelo \textbf{raciocínio} lógico-dedutivo.
Não se trata de trazer para o Ensino Médio a Matemática formal dedutiva, mas de permitir \textbf{que} os jovens percebam a diferença entre \textbf{uma} dedução originária da observação empírica e \textbf{uma} dedução formal.
É importante também verificar \textbf{que} essa competência e suas habilidades não se desenvolvem em separado das demais ; ela é \textbf{um} foco a mais de atenção para o ensino em \textbf{termos} de formação dos estudantes, de modo \textbf{que} identifiquem a Matemática diferenciada das demais Ciências.
As habilidades para essa competência \textbf{demandam} \textbf{que} o estudante vivencie a investigação, a formulação de hipóteses e a tentativa de validação de suas hipóteses.
De certa forma, a proposta é \textbf{que} o estudante do Ensino Médio possa conhecer parte do processo de construção da Matemática, tal qual aconteceu ao longo da história, fruto do pensamento de muitos em diferentes culturas. \textbf{Um} ponto de atenção está no fato de \textbf{que} algumas das habilidades escolhidas pela BNCC ( 2018 ) para essa competência \textbf{remetem} a conteúdos muito específicos, de pouca \textbf{aplicabilidade} e de difícil \textbf{contextualização}, mas \textbf{que}, no entanto, favorecem a investigação e a formulação de hipóteses antes de \textbf{que} os estudantes conheçam os conceitos ou a \textbf{teoria} subjacente a esses conteúdos específicos.
As habilidades propostas para essa competência possuem níveis diferentes de complexidade cognitiva, desde a identificação de \textbf{uma} propriedade até a investigação completa com dedução de \textbf{uma} regra ou procedimento.
Essa competência se relaciona com as Competências Gerais 2, 4, 5 e 7 da BNCC, \textbf{uma} vez \textbf{que} há o incentivo ao exercício da curiosidade intelectual na investigação, neste caso, com maior \textbf{centralidade} no conhecimento matemático.
A linguagem e os recursos digitais são ferramentas básicas e essenciais para facilitar a observação de regularidades, expressar ideias e construir argumentos com base em fatos.
O conjunto das competências específicas da área de Matemática estão articuladas entre si, possibilitando aos estudantes identificarem a importância do conhecimento matemático em seu desenvolvimento pessoal, social e profissional.
SÍNTESE DAS COMPETÊNCIAS Aplicar conhecimentos Construir Argumentos Comunicar em Matemática Formalizar/Demonstrar PROCESSOS PARA APRENDIZAGEM Investigação Construção de Modelos Resolução de Problemas As habilidades selecionadas para a matriz curricular de matemática As habilidades previstas na matriz curricular de Matemática para a etapa do Ensino Médio na BNCC se articulam com as competências gerais e específicas da área e também com as aprendizagens já previstas para o ensino fundamental.
No entanto, como \textbf{explicitado} no documento da Base : > [ … ] Por sua vez, embora cada habilidade esteja associada a determinada competência, isso não significa \textbf{que} ela não contribua para o desenvolvimento de outras. [ … ] as habilidades são apresentadas sem indicação de seriação.
Essa decisão permite flexibilizar a definição anual dos currículos e propostas pedagógicas de cada escola. ( p. 530 ) Essa característica permite escolher as habilidades livres da pressão da seriação e, por \textbf{um} critério de decidir por sua \textbf{centralidade} na formação integral, pela importância dela na constituição da área, pelas conexões \textbf{que} favorecem na área ou entre áreas e, ainda, pela possibilidade de seu desenvolvimento em períodos reduzidos de escolaridade, tais como programas de aceleração da aprendizagem ou cursos de Educação de Jovens e Adultos ( EJA ).
Na sequência, apresentaremos o quadro destinado à Matriz Curricular de Matemática indicada por série, de modo a possibilitar melhor adequação ao trabalho pedagógico dos professores, ao \textbf{que} inclui -se a carga horária prevista.
Além disso, são especificados as Competências Específicas ( CE ) da área de Matemática ( CE01, CE02, CE03, CE04, CE05 ) e as Competências Gerais da BNCC as quais se relacionam ; os objetos do conhecimento e objetivos de aprendizagem para cada habilidade, segundo a competência específica ( CE ) estabelecida.
Para melhor compreensão, ressaltamos \textbf{que} as habilidades são o pressuposto para a indicação dos objetos do conhecimento e objetivos de aprendizagem, o \textbf{que} permitiu elaborar a matriz por série.
Por isso, é possível \textbf{que} habilidades diferentes incorram em objetos do conhecimento semelhantes ou iguais.
No entanto, o \textbf{que} será determinante na \textbf{abordagem} deste objeto do conhecimento pelo professor é, unicamente, a complexidade \textbf{exigida} pela habilidade.
Assim, por exemplo, o \textbf{tópico} de “ juros \textbf{simples} ” está inserido na 1ª e 3ª séries, porém com habilidades e objetivos de aprendizagem distintos, na primeira são observados aspectos relativos à definição e cálculo e na outra, aplicações deste conceito em objetivos mais complexos da matemática financeira ; o mesmo acontece com outros objetos do conhecimento especificados na matriz.
O currículo proposto nesse formato possibilita \textbf{que} o aluno possa ir do conhecimento geral ao conhecimento especializado naturalmente.
Isso se deve a \textbf{um} modelo de aprendizagem contínua \textbf{que} evita \textbf{que} os conceitos caiam facilmente no esquecimento.
Como mencionado na introdução deste documento, o desenvolvimento integral, as dez competências gerais e os princípios \textbf{metodológicos} integradores são a principal forma de integrar o desenvolvimento de ações educacionais de diversos tipos a partir da matriz curricular.
No entanto, é possível observar situações específicas nas quais a Matemática se relaciona com as demais áreas, possibilitando a construção de conhecimentos de modo contextualizado pelos estudantes.
A Matemática pode ser entendida como \textbf{uma} linguagem, sendo parte dos processos de letrar os alunos em diferentes sentidos, formando leitores para múltiplas formas de comunicação, \textbf{aproximando} -se assim da área de Linguagens.
As Artes Gráficas e a Música são outras promissoras formas de integração da Matemática com a área de Linguagens.
Conhecimentos de geometria ampliam a \textbf{percepção} de espaço e das formas de sua representação, fornecendo referenciais para as representações planas nos desenhos, mapas e ferramentas de localização espacial.
Tais representações são úteis a Ciências da Natureza e Ciências Humanas.
As habilidades relacionadas à probabilidade e estatística, em especial aquelas ligadas a coletar, organizar, descrever e analisar dados, realizar inferências e fazer predições com base em amostras de população, são aplicações da Matemática em questões do mundo real \textbf{que} tiveram \textbf{um} crescimento muito grande e se tornaram instrumentos importantes também para as Linguagens e as Ciências Humanas.
A articulação entre Matemática e Ciências da Natureza se dá pela leitura, pela escrita e pela resolução de problemas, nos quais a Matemática \textbf{que} sintetizam as investigações, a generalização de resultados encontrados e a possibilidade de utilizá -los na resolução de problemas semelhantes.
Além disso, projetos temáticos se constituirão cenários propícios à integração da Matemática com as demais áreas de conhecimento.
Outro aspecto a ser considerado na implementação da matriz curricular é o planejamento.
Planejar a utilização das metodologias integradoras de forma mais adequada para o ensino de cada objeto de conhecimento ou habilidade não pode ser feito no momento da aula, \textbf{uma} vez \textbf{que} é \textbf{preciso} considerar os alunos, os recursos disponíveis, o tempo para o ensino e o acompanhamento das aprendizagens esperadas, sendo \textbf{que} a avaliação é \textbf{central} nesse processo.
A avaliação em seu caráter diagnóstico é feita sempre \textbf{que} for necessário ouvir e observar os estudantes, para permitir \textbf{que} expressem ou registrem o \textbf{que} sabem, ou pensam \textbf{que} sabem, sobre \textbf{um} conceito, \textbf{um} problema, \textbf{uma} observação sistemática etc.
A avaliação deve cumprir também seu papel de acompanhamento do processo de aprendizagem, quando o professor observa, registra, solicita \textbf{que} o estudante fale ou represente como pensou, como fez, ou o \textbf{que} não entendeu.
Na perspectiva de \textbf{uma} avaliação formativa, é \textbf{preciso} considerar as limitações dos instrumentos \textbf{prova}, listas de exercícios ou pesquisas sem alinhamento com \textbf{um} projeto, muito utilizados apenas para gerar notas.
Isso significa \textbf{qualificar} mais o instrumento prova, adotando diferentes modalidades para o emprego desse recurso de avaliação, e inserir outras formas de avaliar validadas pelo registro e pela clareza do \textbf{que} se está avaliando. \textbf{Um} elemento da matriz \textbf{que} é muito relevante para o planejamento e a avaliação está traduzido em “ exemplos de objetivos de aprendizagem ”.
Eles indicam aquilo \textbf{que} se espera de aprendizagem para cada conjunto de competências e habilidades e, podem ser parâmetros para acompanhar ensino e aprendizagem, fazer ajustes na caminhada e planejar formas de \textbf{conseguir} \textbf{que} todos aprendam.

% matched lemmas: abordagem, aplicabilidade, aprimorar, aproximar, argumentar, basear, central, centralidade, conseguir, contemporâneo, contextualização, convir, demandar, exigir, explicitar, frisar, globalizado, implicar, mero, metodológico, notar, percepção, posicionamento, preciso, prova, qualificar, que, raciocínio, remeter, restringir, revelar, simples, teoria, termos, tópico, um
\end{textsample}
